\begin{abstract}
We develop Hypothesis H1 from a positive finite transfer amplitude through
vacuum normalization, state-weighted defect transfer, connected-chain bounds,
and the current Yang--Mills analytic frontier.  The development corrects an
important convention: the positive vacuum eigenfunction \(\psi\) enters the
Doob transform, while the natural forward rare-state cost weight is
\(V=\psi^{-1}\).  For a masked bad kernel this choice cancels the vacuum
amplitude ratios exactly.  Independently, the Perron vector of a bad
subkernel minimizes its positive weighted ceiling.  Applying this theorem to
the package's four-channel star kernel proves that the previously selected
weight is exactly optimal: its weighted ceiling is
\(0.002430277501841998\), while its scalar row ceiling is one.  The resulting
degree-18 connected-cluster bound is strictly subcritical.  These statements
close H1 for finite positive kernels and conditional transfer amplitudes.  A
physical Yang--Mills theorem still requires construction of the normalized
nonlinear retained-state kernel, a reflection-compatible eigenweight, and
volume- and cutoff-uniform bad-row and root-moment bounds.
\end{abstract}
